\subsection{Example: Symbolic Differentiation}
\label{Section 2.3.2}

As an illustration of symbol manipulation and a further illustration of data
abstraction, consider the design of a procedure that performs symbolic
differentiation of algebraic expressions.  We would like the procedure to take
as arguments an algebraic expression and a variable and to return the
derivative of the expression with respect to the variable.  For example, if the
arguments to the procedure are \( ax^2 + bx + c \) and \( x \), the procedure
should return \( 2ax + b \).  Symbolic differentiation is of special historical
significance in Lisp.  It was one of the motivating examples behind the
development of a computer language for symbol manipulation.  Furthermore, it
marked the beginning of the line of research that led to the development of
powerful systems for symbolic mathematical work, which are currently being used
by a growing number of applied mathematicians and physicists.

\begin{revealTheSurface}
Python is unusually good at this, and by the end of the section we shall be
writing

\begin{codeBlock}
>>> deriv(x + 3, x)
1
\end{codeBlock}

\noindent
with the expression written the way one would write it on paper.  We are not
going to begin there.  We shall first build the differentiator the way the
original does, with the expression held as data and taken apart by procedures
we write, because that version is the one that shows what the convenient
version is standing on---and because two of this section's exercises exist only
in the first version, and are worth doing before they dissolve.
\end{revealTheSurface}

\subsubsection*{The differentiation program with abstract data}

In order to keep things simple, we will consider a very simple
symbolic-differentiation program that handles expressions that are built up
using only the operations of addition and multiplication with two arguments.
Differentiation of any such expression can be carried out by applying the
following reduction rules:

$$
\frac{{\it dc}}{{\it dx}} = 0,
\quad {\rm for\ } c\ {\rm a\ constant\ or\ a\ variable\ different\ from\ } x,
$$

$$
\frac{{\it dx}}{{\it dx}} = 1,
$$

$$
\frac{{\it d\,(u + v\,)}}{{\it dx}}
=
\frac{{\it du}}{{\it dx}} + \frac{{\it dv}}{{\it dx}},
$$

$$
\frac{{\it d\,(uv\,)}}{{\it dx}}
=
u \frac{{\it dv}}{{\it dx}} + v \frac{{\it du}}{{\it dx}}.
$$

\noindent
Observe that the latter two rules are recursive in nature.  That is, to obtain
the derivative of a sum we first find the derivatives of the terms and add
them.  Each of the terms may in turn be an expression that needs to be
decomposed.  Decomposing into smaller and smaller pieces will eventually
produce pieces that are either constants or variables, whose derivatives will
be either 0 or 1.

To embody these rules in a procedure we indulge in a little wishful thinking,
as we did in designing the rational-number implementation.  If we had a means
for representing algebraic expressions, we should be able to tell whether an
expression is a sum, a product, a constant, or a variable.  We should be able
to extract the parts of an expression.  For a sum, for example we want to be
able to extract the addend (first term) and the augend (second term).  We
should also be able to construct expressions from parts.  Let us assume that we
already have procedures to implement the following selectors, constructors, and
predicates:

\begin{syntaxForm}
is_variable(e)           Is e a variable?
is_same_variable(v1, v2) Are v1 and v2 the same variable?
is_sum(e)                Is e a sum?
addend(e)                Addend of the sum e.
augend(e)                Augend of the sum e.
is_product(e)            Is e a product?
multiplier(e)            Multiplier of the product e.
multiplicand(e)          Multiplicand of the product e.
make_sum(a1, a2)         Construct the sum of a1 and a2.
make_product(m1, m2)     Construct the product of m1 and m2.
\end{syntaxForm}

\noindent
Using these, and a predicate \code{is\_number} which identifies numbers, we can
express the differentiation rules as the following procedure:

\begin{codeBlock}
>>> def deriv(exp, var):
...     if is_number(exp):
...         return 0
...     elif is_variable(exp):
...         return 1 if is_same_variable(exp, var) else 0
...     elif is_sum(exp):
...         return make_sum(deriv(addend(exp), var),
...                         deriv(augend(exp), var))
...     elif is_product(exp):
...         return make_sum(
...             make_product(multiplier(exp),
...                          deriv(multiplicand(exp), var)),
...             make_product(deriv(multiplier(exp), var),
...                          multiplicand(exp)))
...     else:
...         raise ValueError(f"unknown expression type: {exp}")
\end{codeBlock}

\noindent
Note the last clause.  Where the original calls \code{error}, we
\code{raise}---the first time in this book that we have had to signal a
failure.\footnote{\code{raise} interrupts the procedure and hands an object
describing the trouble back to whoever called it.  \code{ValueError} is one of
Python's ready-made kinds of trouble, and the \code{f} before the string marks
it as a \newterm{formatted} string, in which an expression between braces is
replaced by its value.  We shall not need more of exceptions than this until
\link{Chapter 3}.}

This \code{deriv} procedure incorporates the complete differentiation
algorithm.  Since it is expressed in terms of abstract data, it will work no
matter how we choose to represent algebraic expressions, as long as we design a
proper set of selectors and constructors.  This is the issue we must address
next.

\subsubsection*{Representing algebraic expressions}

We can imagine many ways to represent algebraic expressions.  For example, we
could use sequences of symbols that mirror the usual algebraic notation,
representing \( ax + b \) as \code{(a, mul, x, add, b)}.  However, one
especially straightforward choice is to put the operator first, and represent
\( ax + b \) as \code{(add, (mul, a, x), b)}.\footnote{This is where the
original and the translation part company quietly.  The original chooses prefix
notation because it is \emph{Lisp's own}: an expression and the program that
manipulates it are written alike, which is the whole reason quotation was
needed in \link{Section 2.3.1}.  Our tuples enjoy no such correspondence.
\code{(add, (mul, a, x), b)} is a perfectly good representation and has nothing
whatever to do with how Python spells \code{a * x + b}.  We choose it because it
is simple, not because it is the language.}

We shall want two symbols to mark the operators, and here the technique of
\link{Section 2.3.1} earns its place: bind each name to the string of its own
name, and expressions can be written without quotation marks anywhere.

\begin{codeBlock}
>>> add = "add"
>>> mul = "mul"
>>> x = "x"
>>> y = "y"
\end{codeBlock}

\noindent
Then our data representation for the differentiation problem is as follows:

\begin{itemize}
\item
The variables are symbols, which for us are strings:

\begin{codeBlock}
>>> def is_variable(e):
...     return isinstance(e, str)
\end{codeBlock}

\item
Two variables are the same if the symbols representing them are equal:

\begin{codeBlock}
>>> def is_same_variable(v1, v2):
...     return is_variable(v1) and is_variable(v2) and v1 == v2
\end{codeBlock}

\item
Numbers are identified by their type, as sums and products will be:

\begin{codeBlock}
>>> def is_number(e):
...     return isinstance(e, int | float)
\end{codeBlock}

\item
Sums and products are constructed as sequences:

\begin{codeBlock}
>>> def make_sum(a1, a2):
...     return (add, a1, a2)

>>> def make_product(m1, m2):
...     return (mul, m1, m2)
\end{codeBlock}

\item
A sum is a sequence whose first element is the symbol \code{add}:

\begin{codeBlock}
>>> def is_sum(e):
...     return isinstance(e, tuple) and e[0] == add
\end{codeBlock}

\item
The addend is the second item of the sum:

\begin{codeBlock}
>>> def addend(s):
...     return s[1]
\end{codeBlock}

\item
The augend is the third item of the sum:

\begin{codeBlock}
>>> def augend(s):
...     return s[2]
\end{codeBlock}

\item
A product is a sequence whose first element is the symbol \code{mul}:

\begin{codeBlock}
>>> def is_product(e):
...     return isinstance(e, tuple) and e[0] == mul
\end{codeBlock}

\item
The multiplier is the second item of the product:

\begin{codeBlock}
>>> def multiplier(p):
...     return p[1]
\end{codeBlock}

\item
The multiplicand is the third item of the product:

\begin{codeBlock}
>>> def multiplicand(p):
...     return p[2]
\end{codeBlock}
\end{itemize}

\noindent
Thus, we need only combine these with the algorithm as embodied by \code{deriv}
in order to have a working symbolic-differentiation program.  Let us look at
some examples of its behavior:

\begin{codeBlock}
>>> deriv((add, x, 3), x)
('add', 1, 0)
>>> deriv((mul, x, y), x)
('add', ('mul', 'x', 0), ('mul', 1, 'y'))
\end{codeBlock}

\noindent
Notice what the notation has bought.  The expression is \emph{written}
\code{(add, x, 3)}, with no quotation marks, because \code{add} and \code{x}
are names whose values are their own spellings; it is \emph{printed}
\code{('add', 1, 0)}, because printing shows the values.  The original writes
\code{'(+ x 3)} and prints \code{(+ 1 0)}, and the difference between the two
forms is exactly the same difference.

The program produces answers that are correct; however, they are unsimplified.
It is true that

$$ \frac{{\it d\,(xy\,)}}{{\it dx}} = x \cdot 0 + 1 \cdot y, $$

\noindent
but we would like the program to know that \( x \cdot 0 = 0 \), \( 1 \cdot y = y \),
and \( 0 + y = y \).  The answer for the second example should have been
simply \code{y}.  When the expressions are complex this becomes a serious
issue: the derivative of \code{(mul, (mul, x, y), (add, x, 3))} comes back as a
sequence too wide to print on this page.

Our difficulty is much like the one we encountered with the rational-number
implementation: we haven't reduced answers to simplest form.  To accomplish the
rational-number reduction, we needed to change only the constructors and the
selectors of the implementation.  We can adopt a similar strategy here.  We
won't change \code{deriv} at all.  Instead, we will change \code{make\_sum} so
that if both summands are numbers, \code{make\_sum} will add them and return
their sum.  Also, if one of the summands is 0, then \code{make\_sum} will return
the other summand.

\begin{codeBlock}
>>> def make_sum(a1, a2):
...     if is_eq_number(a1, 0):
...         return a2
...     elif is_eq_number(a2, 0):
...         return a1
...     elif is_number(a1) and is_number(a2):
...         return a1 + a2
...     else:
...         return (add, a1, a2)
\end{codeBlock}

\noindent
This uses the procedure \code{is\_eq\_number}, which checks whether an
expression is equal to a given number:

\begin{codeBlock}
>>> def is_eq_number(e, num):
...     return is_number(e) and e == num
\end{codeBlock}

\noindent
Similarly, we will change \code{make\_product} to build in the rules that 0
times anything is 0 and 1 times anything is the thing itself:

\begin{codeBlock}
>>> def make_product(m1, m2):
...     if is_eq_number(m1, 0) or is_eq_number(m2, 0):
...         return 0
...     elif is_eq_number(m1, 1):
...         return m2
...     elif is_eq_number(m2, 1):
...         return m1
...     elif is_number(m1) and is_number(m2):
...         return m1 * m2
...     else:
...         return (mul, m1, m2)
\end{codeBlock}

\noindent
Here is how this version works on our three examples:

\begin{codeBlock}
>>> deriv((add, x, 3), x)
1
>>> deriv((mul, x, y), x)
'y'
>>> deriv((mul, (mul, x, y), (add, x, 3)), x)
('add', ('mul', 'x', 'y'), ('mul', 'y', ('add', 'x', 3)))
\end{codeBlock}

\noindent
\code{deriv} was not touched.  Every one of those improvements was made by
editing two constructors, and that is the whole argument for the abstraction
barrier: the algorithm and the representation were separated far enough that
the second could be improved without disturbing the first.

Although this is quite an improvement, the third example shows that there is
still a long way to go before we get a program that puts expressions into a
form that we might agree is ``simplest.''  The problem of algebraic
simplification is complex because, among other reasons, a form that may be
simplest for one purpose may not be for another.

\begin{exerciseGroup}
\phantomsection\label{Exercise 2.56}
\exerciseHeading{Exercise 2.56:} Show how to extend the basic
differentiator to handle more kinds of expressions.  For instance, implement
the differentiation rule

$$ \frac{{\it d\,(u^n\,)}}{{\it dx}} = nu^{n-1} \frac{{\it du}}{{\it dx}} $$

\noindent
by adding a new clause to the \code{deriv} program and defining appropriate
procedures \code{is\_exponentiation}, \code{base}, \code{exponent}, and
\code{make\_exponentiation}.  (You may bind a symbol \code{power} to denote
exponentiation, as we did for \code{add} and \code{mul}.)  Build in the rules
that anything raised to the power 0 is 1 and anything raised to the power 1 is
the thing itself.

\phantomsection\label{Exercise 2.57}
\exerciseHeading{Exercise 2.57:} Extend the differentiation
program to handle sums and products of arbitrary numbers of (two or more)
terms.  Then the last example above could be expressed as

\begin{codeBlock}
>>> deriv((mul, x, y, (add, x, 3)), x)
\end{codeBlock}

Try to do this by changing only the representation for sums and products,
without changing the \code{deriv} procedure at all.  For example, the
\code{addend} of a sum would be the first term, and the \code{augend} would be
the sum of the rest of the terms.

\phantomsection\label{Exercise 2.58}
\exerciseHeading{Exercise 2.58:} Suppose we want to modify the
differentiation program so that it works with ordinary mathematical notation,
in which \code{add} and \code{mul} are infix rather than prefix operators.
Since the differentiation program is defined in terms of abstract data, we can
modify it to work with different representations of expressions solely by
changing the predicates, selectors, and constructors that define the
representation of the algebraic expressions on which the differentiator is to
operate.

\begin{enumerate}[a.]

\item
Show how to do this in order to differentiate algebraic expressions presented
in infix form, such as \code{(x, add, (3, mul, (x, add, (y, add, 2))))}.  To
simplify the task, assume that \code{add} and \code{mul} always take two
arguments and that expressions are fully parenthesized.

\item
The problem becomes substantially harder if we allow standard algebraic
notation, such as \code{(x, add, 3, mul, (x, add, y, add, 2))}, which drops
unnecessary parentheses and assumes that multiplication is done before
addition.  Can you design appropriate predicates, selectors, and constructors
for this notation such that our derivative program still works?\footnote{Part
(b) is a small precedence parser, and the whole of it turns on one observation:
\code{add} binds more loosely than \code{mul}, so an expression containing a
top-level \code{add} \emph{is} a sum, and may be split at the first one---
everything before it is the addend, everything after is the augend, and the
augend may still contain operators, which the recursion will deal with.  Eight
lines suffice.  Do not be put off by the original's warning that this is
substantially harder; it is harder than part (a), which is six one-line
procedures, but it is not long.}

\end{enumerate}
\end{exerciseGroup}

\subsubsection*{The same program, written in Python}

A word about what follows, because it arrives out of order.  The rest of this
section uses \newterm{classes}, which we have not introduced and shall not
properly introduce until \link{Section 2.4} and \link{Section 2.5}.  Read it
for the shape rather than the detail; the code runs, and the machinery under it
is explained in due course.

The order is not mere convenience.  The two sections that follow are about
precisely the problems a class solves, and the original must solve them by hand
because Scheme offers nothing of the kind.  \link{Section 2.4} is about
attaching a tag to a piece of data so that a program can ask what sort of thing
it is holding, and about dispatching to the right procedure once it knows---%
which is what a class is, and what a method is.  \link{Section 2.5} is about
making one operation serve many kinds of data, which is what \code{\_\_add\_\_}
is for.  Watching the original build both by hand is worth doing, and worth
more once one has seen the finished article.

This is also the second of the two passes promised at the opening of the
chapter: first the original's version, so that its argument stays visible, and
then the version a Python programmer would actually write.  This is the section
where the difference between them is largest, and it would be a poor bargain to
spend it and say nothing.

\begin{revealTheSurface}
We promised at the outset that this would end with \code{deriv(x + 3, x)}, and
we can now say what that costs and what it buys.  Everything above kept the
expression at arm's length: a tuple, tagged with a symbol, taken apart by
selectors we wrote.  Python will let us do something the original cannot, which
is to make \code{x + 3} build the expression itself.

To get there we need \code{class}, which is a third way of defining something,
and worth naming as such.  \code{=} gives a name to a value; \code{def} gives a
name to a procedure; \code{class} gives a name to a \emph{kind} of thing, and
says what procedures come with it.\footnote{Though the three are less separate
than they look.  A procedure made with \code{def} is itself an object of a
class---\code{type(square)} answers \code{<class 'function'>}---and it can be
called at all because that class supplies a \code{\_\_call\_\_}, which is a
two-underscore name of just the sort described below and available to any class
we care to write.  So a procedure is an object that happens to be callable, and
an object of ours becomes callable the moment we give it a \code{\_\_call\_\_}
of its own.  \code{def} is a shortcut for making one, though not quite a
shortcut for \code{\_\_init\_\_}: Python builds a function with an instruction
of its own rather than by calling a constructor.  In \link{Section 2.1.3} we
represented data as a procedure; here the language represents a procedure as
data.  The two directions meet.}

It is closer to what we have already built than it looks.  In
\link{Section 2.1.3} we made a pair out of a procedure that had captured
\code{x} and \code{y}, and we looked at the captured values directly through
\code{\_\_closure\_\_}.  A class is that idea made comfortable.
\code{\_\_init\_\_} is the procedure that does the capturing: it runs when the
object is made, so that writing \code{Symbol("x")} calls it with \code{"x"},
and what it stores in \code{self} is what the other procedures can afterwards
see.  The rest are ordinary procedures that share those captured names.

Names with two underscores at each end, like that one, are the procedures
Python calls on our behalf rather than waiting to be called, and they do very
nearly what one would guess: \code{\_\_add\_\_} for \code{+},
\code{\_\_mul\_\_} for \code{*}, \code{\_\_eq\_\_} for \code{==},
\code{\_\_repr\_\_} for printing.  There are a good many more, and
\link{Section 2.5} makes proper use of them; four will do here.

The first of those is the one that matters to us.  When Python evaluates
\code{a + b} it does not add anything; it asks \code{a} to add \code{b}, by
calling \code{a.\_\_add\_\_(b)}.  So if \code{a} is a kind of object we have
defined, then what \code{+} means for it is ours to say---and \code{x + 3} can
build an expression rather than compute a number.\footnote{And if \code{a} does
not know what to do with \code{b}, Python tries \code{b.\_\_radd\_\_(a)}---the
same operation approached from the right.  That is how \code{3 * x} works as
well as \code{x * 3}, and it does more work here than it may appear to: half
the values \code{deriv} produces are plain numbers, so the addition in its sum
clause is very often \code{0 + y}, which Python resolves by asking \code{y}.}

\begin{codeBlock}
>>> class Expression:
...     def __add__(self, other):
...         return make_sum(self, other)
...     def __radd__(self, other):
...         return make_sum(other, self)
...     def __mul__(self, other):
...         return make_product(self, other)
...     def __rmul__(self, other):
...         return make_product(other, self)
\end{codeBlock}

\noindent
A symbol is an expression that knows its own name; a sum and a product are
expressions that know their parts.  These three are our representation, and
they replace the tuples---nothing else changes.

\begin{codeBlock}
>>> class Symbol(Expression):
...     def __init__(self, name):
...         self.name = name
...     def __repr__(self):
...         return self.name
...     def __eq__(self, other):
...         return isinstance(other, Symbol) and self.name == other.name

>>> class Sum(Expression):
...     def __init__(self, addend, augend):
...         self.addend, self.augend = addend, augend
...     def __repr__(self):
...         return f"({self.addend} + {self.augend})"

>>> class Product(Expression):
...     def __init__(self, multiplier, multiplicand):
...         self.multiplier = multiplier
...         self.multiplicand = multiplicand
...     def __repr__(self):
...         return f"({self.multiplier} * {self.multiplicand})"
\end{codeBlock}

\noindent
The selectors are now attributes---\code{exp.addend} where we wrote
\code{addend(exp)}---and the predicates are \code{isinstance} checks.  But the
constructors are the same two procedures, with the same simplifying rules
inside them, and \code{deriv} is the same algorithm:

\begin{codeBlock}
>>> def deriv(exp, var):
...     if is_number(exp):
...         return 0
...     elif isinstance(exp, Symbol):
...         return 1 if exp == var else 0
...     elif isinstance(exp, Sum):
...         return deriv(exp.addend, var) + deriv(exp.augend, var)
...     elif isinstance(exp, Product):
...         return (exp.multiplier * deriv(exp.multiplicand, var)
...                 + deriv(exp.multiplier, var) * exp.multiplicand)
...     else:
...         raise ValueError(f"unknown expression type: {exp}")
\end{codeBlock}

\noindent
Look closely at the sum clause.  It reads \code{deriv(...) + deriv(...)}---an
ordinary addition---and that \code{+} is \code{\_\_add\_\_}, which is
\code{make\_sum}, which simplifies.  The constructors have not been abolished;
they have been moved behind the operators, so that writing the rule the way a
mathematician writes it invokes them automatically.  The simplification of the
last few pages now happens because of how the rule is spelled.

\begin{codeBlock}
>>> x, y = Symbol("x"), Symbol("y")
>>> x + 3
(x + 3)
>>> deriv(x + 3, x)
1
>>> deriv(x * y, x)
y
>>> deriv(x * y * (x + 3), x)
((x * y) + (y * (x + 3)))
\end{codeBlock}

\noindent
Now see what became of the exercises.  \link{Exercise 2.56} is still real: it
asks for a rule we have not written, and becomes a \code{Power} class and an
\code{\_\_pow\_\_} method.  \link{Exercise 2.57} is still real, and is now about
arity rather than notation---what should \code{Sum} hold if a sum may have any
number of terms?  \link{Exercise 2.58} has gone.  Part (a) asked for infix, and
we have infix.  Part (b), the one the original calls substantially harder,
asked for precedence and for unnecessary parentheses to be dropped, and Python's
own parser has done both before our code is reached.

That last one is worth being precise about, because it is easy to conclude too
much from it.  Python has not given us algebra for nothing.  It has given us
the \emph{parser}---the reading of \code{x + 3 * y} as a product inside a
sum---and that is a real gift, being the eight lines of \link{Exercise 2.58}
and rather more besides.  What it has not given us is any of the meaning.
\code{\_\_add\_\_} and \code{\_\_mul\_\_} we wrote ourselves, and they work only
because we control both kinds of object involved.  The convenience is real, and
it is standing on the same constructors and selectors we spent the first half
of this section building by hand.
\end{revealTheSurface}
